\chapter*{摘要}

\begingroup
\setlength{\baselineskip}{20pt}

燃气-蒸汽联合循环机组具有发电效率高、启停灵活和污染物排放较低等优点，是燃气发电和热电联产系统中的重要装备。对于二拖一燃气-蒸汽联合循环热电联产机组而言，燃气轮机、余热锅炉、汽轮机和供热系统之间存在强耦合关系，运行工况、环境条件和供热需求变化会共同影响机组能量利用水平。因此，建立能够反映部件特性和整机性能的热力分析模型，对机组运行评价具有重要意义。

本文以某二拖一燃气-蒸汽联合循环热电联产机组为研究对象，基于历史运行数据建立了面向工程应用的性能分析方法。首先，构建了统一的工质流股描述方法，对空气、天然气、烟气以及水和水蒸气等工质的温度、压力、质量流量、比焓、比熵和比㶲进行统一封装，为各设备模型之间的数据传递提供基础。其次，建立了燃气轮机部件级模型，将压气机、燃烧室和透平分别作为控制体进行能量分析，然后，建立了余热锅炉和汽轮机模型，从烟气侧放热、水侧吸热、汽轮机各缸做功和供热工况切换等方面分析汽水侧性能。最后，结合 Pearson 相关性分析，研究了非供热工况下联合循环发电效率以及供热工况下热电效率的主要影响因素。

研究结果表明，燃气轮机能量分配结果显示，燃料输入在发电输出和排气余热之间的分配关系对联合循环性能具有重要影响。汽水侧分析表明，余热锅炉高压蒸汽吸热和汽轮机高、中压缸出力对底循环发电能力贡献显著，供热工况下蒸汽能量在发电和供热之间重新分配。相关性分析表明，非供热工况下联合循环发电效率主要受汽机出力、负荷水平、主蒸汽参数和余热回收能力影响；供热工况下热电效率则主要受供热量、环境条件和燃机负荷匹配影响。

本文建立的模型能够在实际运行数据基础上实现燃气轮机、余热锅炉、汽轮机和联合循环整机的性能评价，为热电联产机组运行诊断和效率分析提供了方法参考。

\vspace{1em}

\noindent{\heiti 关键词：}燃气-蒸汽联合循环；热电联产；效率分析；能量分析

\endgroup

\newpage

\begingroup
\titleformat{name=\chapter,numberless}
  {\centering\sffamily\fontsize{16pt}{20pt}\selectfont}
  {}
  {0pt}
  {}
\chapter*{Abstract}
\endgroup

\begingroup
\setlength{\baselineskip}{20pt}

Gas-steam combined cycle units are widely used in power generation and combined heat and power systems because of their high efficiency, operational flexibility and relatively low emissions. For a two-on-one combined cycle cogeneration unit, the gas turbines, heat recovery steam generators, steam turbine and heating system are strongly coupled. Variations in operating load, ambient conditions and heating demand jointly affect the energy utilization of the whole plant. Therefore, a thermodynamic performance analysis model that can describe both component behavior and overall plant performance is important for operation evaluation.

This thesis studies a two-on-one gas-steam combined cycle cogeneration unit and develops an engineering-oriented performance analysis method based on historical operating data. First, a unified working-fluid stream description method is established to represent the temperature, pressure, mass flow rate, specific enthalpy, specific entropy and specific exergy of air, natural gas, flue gas, water and steam. This provides a consistent data interface for different equipment models. Second, a component-level gas turbine model is developed by treating the compressor, combustor and turbine as separate control volumes for energy analysis. Third, models of the heat recovery steam generator and steam turbine are established to analyze flue-gas heat release, water-side heat absorption, cylinder power distribution and heating operation. Finally, Pearson correlation analysis is used to investigate the main factors affecting combined cycle power generation efficiency under non-heating conditions and heat-electricity efficiency under heating conditions.

The results show that the gas turbine energy distribution indicates that the allocation of fuel input between electric power output and exhaust heat has a significant influence on combined cycle performance. The steam-water side analysis shows that high-pressure steam heat absorption and the output of the high- and intermediate-pressure turbine cylinders contribute significantly to bottoming-cycle power generation, while steam energy is redistributed between power generation and heating under heating conditions. Correlation analysis indicates that the combined cycle power generation efficiency under non-heating conditions is mainly affected by steam turbine output, load level, main steam parameters and heat recovery capability, whereas the heat-electricity efficiency under heating conditions is mainly affected by heating load, ambient conditions and gas turbine load matching.

The model developed in this thesis can evaluate the performance of the gas turbine, heat recovery steam generator, steam turbine and overall combined cycle unit based on operating data, providing a methodological reference for operation diagnosis and efficiency analysis of cogeneration units.

\vspace{1em}

\noindent\textbf{Keywords:} gas-steam combined cycle; combined heat and power; efficiency analysis; energy analysis

\endgroup

\newpage
